\section{Experiments}
\label{sec:experimental-design}     % 内部引用符，文章中其他部分可以引用这里

\paragraph{Datasets.}

% 我们的训练数据基于 \citep{sahinuc-etal-2026-reward} 的数据构建，
% 并进行了任务筛选、提示词修改和样本级质量控制。
% 随后，我们采用模型蒸馏为每个样本生成教师 rationale。
% 数据覆盖两类科学写作评测任务：
% 同行评审质量评分（五级）和相关工作评测（二分类）。
% 除标准的 SO 与 RS 训练实例外，
% 我们还基于适配后的样本，
% 为 ReaSon 与 RePaIR 构造了 regenerated-rationale 训练数据。
% PaIR 的配对 DS/RF 训练视图构建方式、
% 完整的数据处理流程及数据统计信息见附录 A。

Our training data is built upon the work of \citet{sahinuc-etal-2026-reward}.
We apply task filtering, prompt revision, and sample-level quality control.
Teacher rationales are then generated via model distillation.
The data covers two scientific writing evaluation tasks:
peer review quality rating (five-point scale) and related-work evaluation (binary classification).
In addition to standard SO and RS training instances,
we construct regenerated-rationale training data
for ReaSon and RePaIR based on the adapted samples.
Details on constructing PaIR's paired DS/RF training views,
the full data processing pipeline, and dataset statistics
are provided in Appendix A.

\paragraph{Models.}

%我们使用的是 Qwen3-4B 模型作为我们的 Base model，所有的训练都在这个模型基础上进行。

We used Qwen3-4B~\citep{yang-etal-2025-qwen3} as the base model for all experiments.
For each task--method combination, we fine-tuned a separate model initialized from 
the same Qwen3-4B checkpoint.

\paragraph{Implementation details.}

%我们使用 LoRA 比较 rationale-and-score 与 score-only 两种监督方式在科学写作评估中的效果。
%七个数据集采用相同的训练、验证和测试划分：训练集用于为每个任务和方法独立训练 adapter，
% 验证集用于逐 epoch 监控，最终保留第 3 个 epoch 的 adapter；测试集用于最终评估。所有评估均采用 zero-shot prompt，模板见图 1。
%LoRA 应用于全部线性层，设置 \(r=16\)、\(\alpha=32\)、学习率 \(1\times10^{-4}\)，训练 3 个 epoch，
% 使用余弦衰减和 8,192-token 最大序列长度。其余配置见附录 B。
%两种监督方式分别通过 direct-score 和 rationale-first 接口进行评估，
%共得到 \(2\times7\times2=28\) 个实验条件，
%每个实验条件均使用 3 个随机种子运行，所有报告结果均为这 3 个随机种子的平均值，
%结果见第~\ref{sec:results} 节。随后，我们分析 RF 接口下的验证集 NLL 和 rationale 质量，
%并对 DS 到 RF 的预测变化进行句子级 rationale 分析。
%对于有害变化，我们替换错误内容并重新预测，以检验 rationale 对分数的影响。最后，分别评估 ReaSon 和 PaIR，
%再验证二者组合 RePaIR 的有效性。

We use LoRA to compare RS and SO supervision for scientific-writing evaluation. 
All seven datasets share the same training, validation, 
and test splits. Separate adapters are trained for each task and method, 
while validation generation is logged after each epoch and the final epoch adapter is retained for final evaluation. 
All evaluations use zero-shot prompts, with templates shown in Figure~\ref{fig:scirm-evaluation}.
LoRA is applied to all linear layers with \(r=16\), \(\alpha=32\), a learning rate of \(1\times10^{-4}\), 
three epochs, cosine decay, and a maximum sequence length of 8,192 tokens. 
Additional settings are provided in Appendix B. 
Evaluating two supervision strategies across seven datasets 
and two inference protocols yields \(2\times7\times2=28\) experimental conditions, 
each of which is run with three random seeds. 
All reported results are averaged across the three seeds (42, 43, and 44) and are 
presented in Section~\ref{sec:results}.
We then analyze validation-set NLL and rationale quality 
under RF, followed by sentence-level analysis of prediction changes from DS to RF. 
For harmful changes, erroneous rationale content is replaced before rescoring to assess its effect on 
the prediction. Finally, we evaluate ReaSon and PaIR separately before testing their combination, RePaIR.

% RQ1：为什么 score-only 训练效果，比 rationale-and-score 训练效果更好
\subsection{Understanding the SO--RS Performance Gap}

% 为检验 SO 与 RS 的性能差距是否与教师 rationale 的拟合程度
% 有关，我们分别计算两个方法的 teacher-forced rationale NLL。
% 分析覆盖所有任务。两个方法采用相同的验证集，共包含 1,368 个
% 样本和 161,178 个教师 rationale token。
%
% 对每个 adapter，逐个输入完整的 prompt--rationale--score 序列，
% 但仅在 rationale token 上计算损失。在因果掩码下，每个 token
% 只能依赖 prompt 和此前的教师 rationale token。随后对全部
% rationale token 进行 micro-average：
% \[
% \operatorname{NLL}_{\mathrm{rat}}
% =
% \frac{
% \sum_n\sum_{t\in R_n}
% -\log p_\theta(r_{n,t}\mid x_n,r_{n,<t})
% }{
% \sum_n |R_n|
% }.
% \]
% 其中，\(R_n\) 表示样本 \(n\) 中的 rationale token 位置。
% 每个方法在每个任务上得到三个 seed-level NLL，并取平均值进行
% 报告。较低的 NLL 表示模型更接近教师 rationale 的分布。

% \paragraph{Teacher-forced rationale NLL 分析。}
\paragraph{Teacher-forced Rationale NLL.}

To test whether the SO--RS performance gap is associated with
teacher-rationale fitting, we compute teacher-forced rationale NLL
for both methods. The analysis covers all tasks. Both methods use identical 
validation splits, comprising 1,368 samples and 161,178 teacher-rationale 
tokens in total.

For each adapter, complete prompt--rationale--score sequences are
processed individually, while loss is computed only at rationale-token
positions. Under causal masking, each token conditions only on the
prompt and preceding teacher rationale tokens. We then micro-average
the loss over all rationale tokens:
\[
\operatorname{NLL}_{\mathrm{rat}}
=
\frac{
\sum_n\sum_{t\in R_n}
-\log p_\theta(r_{n,t}\mid x_n,r_{n,<t})
}{
\sum_n |R_n|
}.
\]
Here, \(R_n\) denotes the rationale-token positions in sample \(n\).
This yields three seed-level NLL values per method and task, which are
averaged for reporting. A lower NLL indicates closer fitting to the
teacher-rationale distribution.

\paragraph{Rationale Quality Comparison.}

% 由于 NLL 只能衡量模型对教师 rationale 的拟合程度，我们进一步
% 使用外部 LLM 评估 SO 和 RS 生成的 rationale 质量及其与预测
% 分数的一致性。该分析固定使用 seed 42 的 RF 输出，覆盖上述
% 七个任务。每个任务随机抽取 50 组 SO--RS 配对样本，共得到
% 350 组样本和 700 条 rationale。
%
% GLM-5.3-Flash 和 Doubao-Seed-2.0-Lite 作为主裁判，
% MiniMax-M3 作为仲裁裁判。当主裁判调用或解析失败、判断不一致，
% 或任一维度相差至少 2 分时，启动仲裁；350 组样本中共有 192 组
% 使用了仲裁。第一轮从评分标准理解、证据支撑、关键信息覆盖和
% 无依据陈述控制四个维度进行 1--3 分评分，并比较平均分。第二轮
% 判断 rationale 是否支持预测分数，结果分为
% \texttt{supported}、\texttt{partially\_supported} 和
% \texttt{not\_supported}。模型配置和评估 prompt 见附录 B。

Because NLL measures only teacher-rationale fitting, we further use
external LLMs to assess generated-rationale quality and consistency
with predicted scores. This analysis uses RF outputs from seed 42
and covers the same seven tasks. We randomly sample 50 paired SO--RS
cases per task, yielding 350 pairs and 700 rationales.

GLM-5.3-Flash and Doubao-Seed-2.0-Lite serve as primary judges, with
MiniMax-M3 as the adjudicator. Adjudication is triggered by invocation
or parsing failure, judge disagreement, or a difference of at least
two points on any dimension. It is used for 192 of the 350 pairs.
The first round scores each rationale from 1 to 3 on rubric
understanding, evidential support, key-information coverage, and
unsupported-claim control. Mean scores determine whether one rationale
is better or both are comparable. The second round classifies whether
each rationale supports its predicted score as \texttt{supported},
\texttt{partially\_supported}, or \texttt{not\_supported}. Model
configurations and evaluation prompts are provided in Appendix B.

% RQ2：为什么在相同训练方法下，RF 推理在序数任务上的 QWK 低于 DS 推理？
\subsection{Understanding the DS--RF Performance Gap on Ordinal Tasks}

% 28 个实验条件表明，推理协议的影响随任务类型和训练目标而变化：
% 序数任务中 DS 均优于 RF；二分类任务中，SO 下 DS 更优，
% 而 RS 下 RF 更优。为分析这一差异，我们比较了两种训练条件下
% DS 到 RF 的样本级预测转移，以及 Accuracy、QWK 和各类结果的变化。
% 该分析覆盖四个序数任务和三个随机种子，每种训练方式包含
% 11,364 组 DS--RF 配对预测。序数预测分为正确、相邻错误、
% 严重错误和格式无效。
%
% 对预测发生变化的样本，我们沿用三裁判机制。裁判综合测试样本、
% gold score、DS 与 RF 预测及 RF rationale，判断 rationale
% 支持 gold score、错误预测，还是缺乏充分证据。对于有害转移，
% 进一步逐句标注 \texttt{factual_error}、
% \texttt{evidence_misread} 等错误类型，完整分类见附录 B。
%
% 最后，我们在 158 个有害转移案例上检验修正 rationale 是否会
% 改变后续评分。原始 rationale 由 GLM-5.3-Flash 修正，并经
% MiniMax-M3 审核；所有案例均使用原 adapter 以及相同的 RF
% system 和 user prompt 进行重评分。我们将自然 RF 生成记为
% NatRF，将固定原错误 rationale 后仅生成分数记为 FixOrig，
% 将固定修正 rationale 后仅生成分数记为 FixCorr。
% NatRF--FixOrig 比较用于验证固定前缀重评分能否复现原预测，
% FixOrig--FixCorr 比较则用于衡量 rationale 修正带来的分数变化。

Across the 28 experimental conditions, the effect of the inference
protocol varied by task type and training objective. DS consistently
outperformed RF on ordinal tasks. On binary tasks, DS performed better
under SO, whereas RF performed better under RS. We therefore compared
sample-level DS-to-RF transitions, together with changes in Accuracy,
QWK, and outcome distributions. The analysis covered four ordinal tasks
and three random seeds, yielding 11,364 paired DS--RF predictions per
training objective. Ordinal predictions were categorized as correct,
adjacent errors, severe errors, or format-invalid.

For samples whose predictions changed, we applied the same three-judge
protocol. Given the test instance, gold score, DS and RF predictions,
and RF rationale, the judges determined whether the rationale supported
the gold score, the incorrect prediction, or neither due to insufficient
evidence. Harmful transitions were further annotated sentence by sentence for
errors such as \texttt{factual\_error} and
\texttt{evidence\_misread}. The complete taxonomy is provided in
Appendix A.

Finally, we tested whether rationale correction changed subsequent
scores in 158 harmful-transition cases. GLM-5.3-Flash corrected the
original rationales, which were then reviewed by MiniMax-M3. Each case
was rescored using its original adapter and the same RF system and user
prompts. We denote natural RF generation as NatRF, rescoring from the
fixed original rationale as FixOrig, and rescoring from the fixed
corrected rationale as FixCorr. The NatRF--FixOrig comparison validates
whether fixed-prefix rescoring reproduces the original prediction,
whereas FixOrig--FixCorr measures the score changes associated with
rationale correction.
